\documentclass[12pt]{article}
\usepackage[utf8]{inputenc}
\usepackage[T1]{fontenc}
\usepackage{hyperref}
\title{Mid-Eastern Tharu Language Resources for K-12 Curriculum Development}
\author{Compiled by Claude Sonnet 5 (claude-sonnet-5) for Andr\'as Kornai}
\date{}
\begin{document}
\maketitle

\section{Introduction}
Mid-Eastern Tharu (ISO 639-3 thq, Glottocode koch1247) is an Indo-Aryan language of Nepal with about 340,000 speakers. This report is based on a small number of searches. I found no resource that is specific to Mid-Eastern (Kochila) Tharu. What I did find concerns neighbouring Tharu varieties, especially Dangaura Tharu, and may not transfer to Mid-Eastern Tharu.

\section{Primary Educational Resources}
\begin{itemize}
\item A UCL-hosted chapter describes a Kapilvastu school whose primary textbooks print Dangaura Tharu and Nepali side by side. It also describes a multilingual textbook used in Grades I--III. These are Dangaura materials. \url{https://discovery.ucl.ac.uk/id/eprint/10169436/}
\item World Vision Nepal developed 15 Tharu-language early-grade readers. They are aimed at Rana Tharu communities. \url{https://www.wvi.org/stories/nepal/embracing-teaching-learning-materials-tharu-language}
\end{itemize}

\section{Linguistic and Reference Materials}
A Tribhuvan University thesis on Dangaura Tharu covers mother-tongue proficiency, domains of use and vitality, and lists community language resources. I did not read it in full. \url{https://portal.tu.edu.np/downloads/Dangaura_Tharu_2023_09_20_20_23_06.pdf}

\section{Religious and Community Resources}
None were found in the search results.

\section{Digital and Audio Resources}
None specific to Mid-Eastern Tharu were found.

\section{Curriculum Development Context}
Nepal's constitution and the National Education Policy 2019 support mother-tongue education. The Curriculum Development Centre is the authoritative source for current Tharu textbook availability, and it should be checked directly. Mid-Eastern Tharu materials appear to be scarce or absent online.
\end{document}